% =============================================================================
% Sections 6, 7, 8 of zkMetal paper
% Target: USENIX Security format (2-column, 10pt)
% =============================================================================

\section{Performance Characterization}
\label{sec:perf}

We evaluate zkMetal against four baselines: ICICLE-Metal~\cite{icicle} (the only other Metal ZK library), ICICLE CPU~\cite{icicle} (Ingonyama's CPU backend on the same hardware), Arkworks~\cite{arkworks} (Rust CPU reference), and ICICLE-CUDA~\cite{icicle} on an NVIDIA RTX~3090~Ti (representing the discrete GPU ceiling).
All zkMetal benchmarks run on an Apple M3~Pro (6P+6E~cores, 150~GB/s memory bandwidth, $\sim$3.6~TFLOPS GPU compute).

\subsection{Comparison with Baselines}
\label{sec:baselines}

\paragraph{MSM.}
Table~\ref{tab:msm-comparison} compares MSM implementations across six systems.
zkMetal achieves 27\,ms at $2^{16}$ points and 45\,ms at $2^{18}$, a 33$\times$ improvement over ICICLE-Metal (1475\,ms) and 5.9$\times$ over Arkworks CPU (266\,ms) at $2^{18}$.
ICICLE-Metal v3.8 incurs $\sim$600\,ms constant overhead per call due to license-server communication, but even accounting for this, zkMetal's kernel execution is substantially faster.
MoPro~v2~\cite{mopro}, measured on M3~Air by its developers, reports 678\,ms at $2^{18}$---15$\times$ slower than zkMetal on comparable hardware.

\begin{table}[t]
\centering
\small
\caption{BN254 G1 MSM comparison across implementations.}
\label{tab:msm-comparison}
\begin{tabular}{lrrrrrr}
\toprule
Points & \textbf{zkMetal} & ICICLE- & ICICLE & MoPro & Ark- & ICICLE \\
       & (M3 Pro) & Metal$^1$ & CPU$^1$ & v2$^2$ & works$^2$ & CUDA$^3$ \\
\midrule
$2^{16}$ & \textbf{27\,ms} & 1083\,ms & 114\,ms & 253\,ms & 69\,ms & $\sim$9\,ms \\
$2^{18}$ & \textbf{45\,ms} & 1475\,ms & 556\,ms & 678\,ms & 266\,ms & -- \\
$2^{20}$ & \textbf{119\,ms} & 2590\,ms & 2349\,ms & 1702\,ms & 592\,ms & -- \\
\bottomrule
\end{tabular}
\\[2pt]
\raggedright\scriptsize
$^1$Measured locally. ICICLE-Metal v3.8 has $\sim$600\,ms per-call overhead.\\
$^2$Reported by MoPro blog~\cite{mopro}; M3~Air, different methodology.\\
$^3$RTX~3090~Ti with native 64-bit integer multiply~\cite{icicle}.
\end{table}

The CUDA gap at $2^{16}$ ($\sim$9\,ms vs.\ 27\,ms, a 3$\times$ difference) reflects two hardware disadvantages: the M3~Pro's 150~GB/s memory bandwidth versus 1008~GB/s on the RTX~3090~Ti (6.7$\times$), and the lack of native 64-bit integer multiply (each \texttt{mulhi(ulong,ulong)} decomposes into 4$\times$32-bit MADs).
This gap is \emph{narrower} than the raw bandwidth ratio suggests because unified memory eliminates PCIe transfer overhead and the Pippenger algorithm is latency-sensitive (random-access scatter), not purely bandwidth-limited.

\paragraph{NTT.}
Table~\ref{tab:ntt-comparison} shows NTT performance across fields and sizes.
zkMetal is 30--90$\times$ faster than ICICLE-Metal on BN254 and 90--500$\times$ faster on BabyBear, with the gap widening at larger sizes.
ICICLE-Metal's $\sim$85\,ms per-call overhead is a constant factor, but even subtracting it leaves zkMetal 5--30$\times$ faster on BN254 at $2^{22}$+.

\begin{table}[t]
\centering
\small
\caption{NTT comparison: zkMetal vs.\ ICICLE-Metal (M3~Pro).}
\label{tab:ntt-comparison}
\begin{tabular}{lrrrr}
\toprule
Size & zkMetal & ICICLE & zkMetal & ICICLE \\
     & BN254   & BN254  & BabyBear & BabyBear \\
\midrule
$2^{16}$ & \textbf{0.76\,ms} & 89\,ms & \textbf{0.18\,ms} & 86\,ms \\
$2^{18}$ & \textbf{1.6\,ms}  & 108\,ms & \textbf{0.26\,ms} & 92\,ms \\
$2^{20}$ & \textbf{6.1\,ms}  & 194\,ms & \textbf{0.95\,ms} & 108\,ms \\
$2^{22}$ & \textbf{26\,ms}   & 915\,ms & \textbf{2.8\,ms}  & 181\,ms \\
$2^{24}$ & \textbf{116\,ms}  & 3892\,ms & \textbf{2.0\,ms} & 709\,ms \\
\bottomrule
\end{tabular}
\end{table}

BabyBear at $2^{24}$ achieves 8.5\,B~elements/sec---a regime where Metal's 32-bit ALU is no longer a disadvantage, since each butterfly is a single 32-bit MAD versus the 64~MADs required for a BN254 butterfly.

\paragraph{Proof systems.}
Table~\ref{tab:proof-systems} presents end-to-end proof system benchmarks.
Circle STARK proves at 21\,ms for a $2^{14}$ trace (full GPU pipeline: Circle NTT LDE 0.9\,ms, Keccak Merkle 6\,ms, constraint evaluation 2\,ms, CPU FRI fold 12\,ms).
HyperNova achieves 0.09\,ms per fold, a 40$\times$ improvement from C~CIOS acceleration, Keccak transcript optimization, and pre-computed affine points.
Groth16 proves at 14\,ms for 256~constraints, a 107$\times$ improvement over the unoptimized baseline.

\begin{table}[t]
\centering
\small
\caption{End-to-end proof system benchmarks.}
\label{tab:proof-systems}
\begin{tabular}{lrrr}
\toprule
System & Size & Prove & Verify \\
\midrule
Circle STARK & $2^{14}$ trace & 21\,ms & 28\,ms \\
Plonk (KZG) & 1024 gates & 49\,ms & 2\,ms \\
Groth16 & 256 constraints & 14\,ms & 73\,ms \\
GKR & $2^{10}$ w, d=4 & 7.7\,ms & 8.4\,ms \\
HyperNova & per fold & 0.09\,ms & -- \\
Spartan & $2^{14}$ & 121\,ms & 8\,ms \\
Lasso & $2^{18}$ & 56\,ms & 31\,ms \\
\bottomrule
\end{tabular}
\end{table}

\paragraph{Cross-curve MSM.}
Table~\ref{tab:cross-curve-msm} characterizes MSM across seven curves at $2^{18}$.
BN254 is fastest (45\,ms) due to combined GLV endomorphism and signed-digit optimization.
BLS12-377 (205\,ms) pays a 4.6$\times$ penalty from wider limbs (12$\times$32-bit vs.\ 8$\times$32-bit for BN254).
secp256k1 (113\,ms) cannot use GLV due to the scheduling pathology (\S\ref{sec:pathology}) and relies on signed-digit only.
Pallas and Vesta, as a cycle pair, achieve nearly identical performance (66 and 65\,ms), validating their suitability for recursive proof composition.

\begin{table}[t]
\centering
\small
\caption{Cross-curve MSM at $2^{18}$ points (GPU).}
\label{tab:cross-curve-msm}
\begin{tabular}{lrrr}
\toprule
Curve & Limbs & $2^{16}$ & $2^{18}$ \\
\midrule
BN254 & 8$\times$32 & 27\,ms & 45\,ms \\
Pallas & 8$\times$32 & 39\,ms & 66\,ms \\
Vesta & 8$\times$32 & 39\,ms & 65\,ms \\
secp256k1 & 8$\times$32 & 24\,ms & 113\,ms \\
BLS12-377 & 12$\times$32 & 176\,ms & 205\,ms \\
\bottomrule
\end{tabular}
\end{table}

\subsection{Theoretical Floor Analysis}
\label{sec:floors}

To characterize optimization maturity, we compute hardware-theoretical performance floors for 33~primitives and compare against measured performance.

\paragraph{Methodology.}
For each primitive, we compute three lower bounds:
(1)~\emph{Compute-bound}: total arithmetic operations divided by the GPU's 3.6~TFLOPS peak (counting BN254 \texttt{fp\_mul} as $\sim$64 32-bit multiply-add operations);
(2)~\emph{Memory-bound}: total bytes read and written divided by 150~GB/s peak bandwidth;
(3)~\emph{Dispatch-bound}: number of Metal command buffer dispatches multiplied by the empirical $\sim$0.5\,ms per-dispatch overhead.
The theoretical floor is $\max(\text{compute}, \text{memory}, \text{dispatch})$.
For compound protocols (proof systems), we sum the component floors.
This methodology deliberately provides \emph{optimistic} floors---real performance must also contend with cache miss penalties, register pressure, and occupancy limits.

\paragraph{Results.}
Table~\ref{tab:floors} presents the 33-primitive analysis ranked by headroom (ratio of measured to theoretical floor).
We identify three regimes:

\emph{Near-optimal ($<$2$\times$ headroom):}
Seven primitives operate within 2$\times$ of their theoretical floors.
BabyBear NTT at $2^{24}$ (1.2$\times$) is essentially memory-bandwidth-limited: 16.8\,M elements $\times$ 4~bytes at 150~GB/s yields a 1.7\,ms floor against the measured 2.0\,ms.
HyperNova (1.3$\times$), Circle NTT (1.3$\times$), Groth16 (1.4$\times$), and KZG commit (1.4$\times$) are similarly mature.
These primitives share a common trait: regular, predictable memory access patterns or small element sizes that pack efficiently into cache lines.

\emph{Moderate headroom (2--5$\times$):}
Thirteen primitives have 2--5$\times$ headroom.
Radix sort (2$\times$), Circle STARK (2$\times$), and Lasso (2$\times$) are approaching practical limits.
ECDSA (3$\times$) and Plonk (3$\times$) have room for improvement through further C~CIOS acceleration of hot loops.
Poseidon2 batch (4.5$\times$) is limited by the 22 sequential rounds inherent to the hash construction, which cap parallelism.

\emph{Significant headroom (5--10$\times$):}
Five primitives have 5$\times$ or more headroom.
MSM BN254 ($\sim$9$\times$) and NTT BN254 ($\sim$9$\times$) are the largest-headroom GPU primitives.
The MSM gap is dominated by the random-access scatter pattern in bucket accumulation: each thread writes to a bucket determined by the scalar digit, causing cache-hostile access with no spatial locality.
The NTT gap reflects strided column access in the four-step FFT: 32-byte elements at 128-byte cache line granularity yield only 25\% cache utilization during column passes.

\begin{table}[t]
\centering
\small
\caption{Theoretical floor analysis for 33 primitives (selected). Full table in Appendix~A.}
\label{tab:floors}
\begin{tabular}{llrrl}
\toprule
Primitive & Current & Floor & Head- & Bottleneck \\
          &         &       & room  &            \\
\midrule
\multicolumn{5}{l}{\emph{Near-optimal ($<$2$\times$)}} \\
BabyBear NTT $2^{24}$ & 2.0\,ms & 1.7\,ms & 1.2$\times$ & Bandwidth \\
Circle NTT $2^{20}$ & 4.0\,ms & 3.0\,ms & 1.3$\times$ & Compute $\approx$ BW \\
HyperNova fold & 0.09\,ms & 0.07\,ms & 1.3$\times$ & C CIOS floor \\
KZG commit $2^{10}$ & 0.7\,ms & 0.5\,ms & 1.4$\times$ & Cached affine \\
Groth16 prove 256 & 14\,ms & 10\,ms & 1.4$\times$ & MSM dominated \\
GKR $2^{10}$ d=4 & 7.7\,ms & 5\,ms & 1.5$\times$ & C CIOS + topology \\
\midrule
\multicolumn{5}{l}{\emph{Moderate (2--5$\times$)}} \\
Lasso $2^{18}$ & 56\,ms & 30\,ms & 2$\times$ & Fused GPU \\
Circle STARK $2^{14}$ & 21\,ms & 10\,ms & 2$\times$ & Batched CBs \\
Radix sort $2^{20}$ & 2.1\,ms & 1\,ms & 2$\times$ & Histogram \\
ECDSA batch 64 & 1.7\,ms & 0.5\,ms & 3$\times$ & C CIOS Fr \\
Plonk 1024 & 49\,ms & 15\,ms & 3$\times$ & Poly ops \\
secp256k1 MSM $2^{18}$ & 113\,ms & 30\,ms & 4$\times$ & No GLV \\
Poseidon2 $2^{16}$ & 8.1\,ms & 1.8\,ms & 4.5$\times$ & Sequential rounds \\
\midrule
\multicolumn{5}{l}{\emph{Significant ($>$5$\times$)}} \\
BLS12-377 MSM $2^{18}$ & 205\,ms & 35\,ms & 6$\times$ & 12-limb Fq \\
Basefold $2^{18}$ & 144\,ms & 20\,ms & 7$\times$ & 18 fold rounds \\
FRI fold $2^{20}$ & 1.96\,ms & 0.3\,ms & 7$\times$ & Bandwidth \\
NTT BN254 $2^{22}$ & 26\,ms & 2.9\,ms & 9$\times$ & Strided BW \\
MSM BN254 $2^{18}$ & 45\,ms & 5\,ms & 9$\times$ & Scatter BW \\
\bottomrule
\end{tabular}
\end{table}

\paragraph{Discussion.}
\figurename~\ref{fig:floor-scatter} plots measured performance against theoretical floors for all 33~primitives.
The dominant factor separating near-optimal from high-headroom primitives is \emph{memory access pattern regularity}, not raw compute efficiency.
BabyBear NTT (1.2$\times$) streams 4-byte elements sequentially; MSM BN254 (9$\times$) scatters 96-byte projective points randomly.
Both are running the same GPU, but MSM's effective memory bandwidth is $\sim$11$\times$ lower due to cache-line waste.

This analysis also explains why small-field primitives consistently achieve lower headroom: 4-byte BabyBear elements pack 32 per cache line versus 4 for 32-byte BN254 elements, yielding 8$\times$ better cache utilization for identical access patterns.
The ``small-field thesis''---that modern proof systems migrating to smaller fields play to Apple Silicon's strengths---is thus doubly supported: smaller fields benefit from both the 32-bit ALU match \emph{and} superior cache behavior.

\begin{figure}[t]
\centering
% Placeholder for Figure 7: scatter plot of current perf vs theoretical floor
\fbox{\parbox{0.9\columnwidth}{\centering\vspace{2cm}
\textbf{Figure~\ref{fig:floor-scatter}}: Measured performance vs.\ theoretical floor for 33~primitives. Diagonal line indicates 1$\times$ (optimal). Primitives above the line have optimization headroom; distance from diagonal indicates gap. Near-optimal cluster (BabyBear NTT, HyperNova, KZG) at bottom-left; high-headroom outliers (BN254 MSM, BN254 NTT) at top-right.
\vspace{2cm}}}
\caption{Current performance vs.\ hardware-theoretical floor.}
\label{fig:floor-scatter}
\end{figure}


% =============================================================================

\section{Negative Results and Dead Ends}
\label{sec:dead-ends}

We catalog 30+ optimization attempts that failed, organized by failure mode.
Negative results are underreported in systems research despite their value in preventing duplicated effort and revealing hardware characteristics.
Each entry describes what was attempted, the observed result, and the root cause.

\subsection{Hardware Pathology}
\label{sec:dead-hw}

\paragraph{64-bit limb CIOS on Metal GPU.}
Hypothesis: representing 256-bit field elements as 4$\times$64-bit limbs instead of 8$\times$32-bit should halve the number of multiply-accumulate operations.
Result: 67\% \emph{slower}.
Root cause: Metal GPUs have no native 64-bit integer multiply hardware.
The instruction \texttt{mulhi(ulong, ulong)} decomposes internally into four 32-bit MAD operations~\cite{metal-spec}, yielding zero computational benefit while doubling register pressure (64-bit registers consume two 32-bit slots).
The increased register pressure reduces occupancy, compounding the already-null benefit.
\emph{Lesson:} Metal's 32-bit ALU is not merely a ``preferred'' width---it is the \emph{only} integer multiply width.
Any algorithm that assumes 64-bit integer multiply is free (common in CUDA ZK libraries~\cite{cuzk, zprize2022}) must be redesigned for Metal.

\paragraph{Cooperative SIMD reduce for MSM buckets (w=13).}
Hypothesis: use Metal SIMD shuffle operations to cooperatively reduce adjacent buckets within a SIMD group, reducing global memory traffic.
Result: 2.4$\times$ slower.
Root cause: each projective elliptic curve point is 96~bytes (3$\times$Fp at 32~bytes each).
Metal SIMD shuffles are designed for 4-byte scalar types; shuffling 96~bytes requires 24 separate shuffle operations per point exchange, with each shuffle consuming a full SIMD cycle.
The shuffle overhead exceeds the memory traffic reduction.
\emph{Lesson:} GPU SIMD cooperative patterns are effective only when the exchanged data fits within the register file's natural granularity.
For large algebraic types (field elements, curve points), thread-private accumulation with global atomic-free scatter is preferable.

\paragraph{Custom AArch64 assembly for Montgomery multiply.}
Hypothesis: hand-tuned assembly using ARM64 \texttt{MUL}/\texttt{UMULH} dual-issue scheduling should outperform compiler-generated code.
Result: 23\% slower than LLVM's output for the equivalent C with \texttt{\_\_uint128\_t}.
Root cause: Clang already emits optimal \texttt{MUL}+\texttt{UMULH} pairs for 128-bit multiply and correctly schedules them for the M-series microarchitecture.
This finding was independently confirmed by two other Apple Silicon ZK projects~\cite{lightningjolt, aztecfasttech}.
\emph{Lesson:} on ARM64 with a modern compiler, hand-written Montgomery assembly is counterproductive.
The compiler's register allocator and instruction scheduler have sufficient information to produce near-optimal code for \texttt{\_\_uint128\_t} CIOS patterns.

\paragraph{Fused reduce+bucket\_sum MSM kernel.}
Hypothesis: combining bucket accumulation and bucket summation into a single kernel dispatch saves one command buffer overhead ($\sim$0.5\,ms) and allows data to remain in registers.
Result: correct but 21\% slower.
Root cause: the fused kernel launches only 4096~threads (one per bucket at window width 12), leaving the GPU severely underutilized.
Metal's scheduler cannot efficiently distribute work when the grid is small relative to the GPU's execution unit count, even when per-thread work is substantial.
\emph{Lesson:} on Metal, maintaining high grid occupancy (many threads with moderate work) outperforms low grid occupancy (few threads with heavy work), unlike CUDA where persistent kernel patterns are effective.

\subsection{Algorithmic Dead Ends}
\label{sec:dead-algo}

\paragraph{GLV endomorphism on non-BN254 curves.}
The GLV endomorphism~\cite{glv} halves the scalar width by decomposing each scalar multiplication into two half-width multiplications using the curve's efficiently computable endomorphism.
On BN254, this provides a reliable 22\% improvement because the curve's endomorphism maps cleanly to the 8-limb arithmetic, and the resulting bucket count avoids scheduling pathologies.
On secp256k1, GLV \emph{regresses by 12$\times$} (13569\,ms vs.\ 1133\,ms at $2^{18}$).
The mechanism: GLV doubles the effective number of points (each scalar becomes two half-width scalars), pushing the bucket count to $\sim$32K---directly into the GPU scheduling pathology band (\S\ref{sec:pathology}).
On BLS12-377, GLV regresses by 1.27$\times$ (1155\,ms vs.\ 911\,ms) because the 12-limb Fq field makes each point addition expensive enough that the cost of processing twice as many points exceeds the savings from shorter scalars.
\emph{Lesson:} GLV is profitable only when (a)~the resulting bucket count avoids hardware pathologies and (b)~point addition cost is low relative to the scalar width reduction.
Both conditions are curve- and hardware-specific.

\paragraph{Local-binning radix sort.}
We tested five variants of local-binning optimizations for the GPU radix sort used in MSM: serial ranking, SIMD ranking, linear binary-search scatter, atomic-based scatter, and sequential SIMD claiming.
All performed the same as or worse than the baseline global-histogram sort.
Root cause: Apple Silicon's large unified L2 cache reduces actual write amplification to 6--9$\times$ (versus the 32$\times$ theoretical worst case for random scatter), making the coalescing optimization that motivates local binning on discrete GPUs unprofitable.
The staging buffer overhead, local histogram loading, prefix sums, and per-element binary search collectively exceed the marginal coalescing benefit.
\emph{Lesson:} unified memory with large caches changes the cost model for memory coalescing optimizations.
Techniques developed for discrete GPUs with separate VRAM should be re-evaluated on unified architectures.

\paragraph{Lazy field reduction in Poseidon2.}
Hypothesis: defer modular reduction across multiple Poseidon2 rounds, batching reductions to amortize their cost.
Result: infeasible.
Root cause: after just two Poseidon2 rounds (S-box + MDS matrix), intermediate values exceed $2^{256}$, overflowing the 8-limb representation.
The Poseidon2 MDS matrix has coefficients that amplify value ranges too aggressively for deferred reduction to be practical.
\emph{Lesson:} lazy reduction is viable only when the algebraic structure guarantees bounded growth across the deferred interval.
Hash functions with dense linear layers are poor candidates.

\paragraph{Fused Merkle subtree kernels.}
Two attempts: Poseidon2 and Blake3 backends.
Poseidon2 fused subtree: the tree narrows at upper levels, leaving 80\% of threads idle after the first two levels.
Blake3 fused subtree: regressed 2$\times$ because Blake3's parent compression is too lightweight for shared-memory barrier overhead to break even.
In both cases, the level-by-level approach (one dispatch per tree level) outperforms the fused approach because each dispatch fully utilizes the GPU at its level.
\emph{Lesson:} fused tree kernels are profitable only when the per-node computation is expensive enough to amortize synchronization barriers \emph{and} the tree is wide enough to maintain occupancy across levels.

\paragraph{BN254 NTT coalescing at $2^{22}$.}
Three approaches to improving the strided column access pattern in four-step FFT:
(1)~Transpose-before-column: four transposes + four blits cost $\sim$50\,ms, exceeding any coalescing benefit for 32-byte elements (482\,ms vs.\ 351\,ms baseline).
(2)~Swapped split ($N_1=10$, $N_2=12$ instead of $N_1=11$, $N_2=11$): 420\,ms with correctness failures---column cache waste is symmetric regardless of split direction.
(3)~Separate command buffers for phase profiling: GPU hang due to buffer ownership conflicts across command buffers with shared scratch memory.
The fundamental issue is that 32-byte elements at 128-byte cache lines yield 25\% utilization during column passes, and no rearrangement of the four-step algorithm can avoid this without introducing costs that exceed the benefit.
\emph{Lesson:} for wide field elements, the NTT's strided access pattern is a fundamental hardware--algorithm mismatch, not an implementation deficiency.

\paragraph{Circle STARK traceLDE readback elimination.}
Hypothesis: keep trace LDE buffers on GPU (avoiding CPU readback) since they are consumed by subsequent GPU Merkle and constraint evaluation kernels.
Result: 38$\times$ regression (21\,ms $\to$ 800\,ms).
Root cause: the CPU FRI fold phase requires the LDE data on the CPU side, and without the readback, subsequent operations stall or consume incorrect data.
The readback itself costs only $\sim$0.1\,ms---eliminating it saves negligible time while breaking the pipeline.
\emph{Lesson:} in unified memory architectures, readback costs are already minimal; the optimization target should be the \emph{computation}, not the data movement.

\subsection{Metal Compiler and Runtime Issues}
\label{sec:dead-compiler}

\paragraph{Inline miscompilation of large-struct functions.}
Metal shader functions returning \texttt{PointProjective} (96~bytes: three Fp values of 32~bytes each) produce incorrect results when the compiler inlines them.
Discovered when 5/9 MSM tests failed after removing \texttt{\_\_attribute\_\_((noinline))} from \texttt{point\_double}, \texttt{point\_add\_mixed}, and \texttt{point\_add}.
The bug is deterministic and reproducible: the same inputs produce correct results with \texttt{noinline} and incorrect results without it.
We attribute this to a register allocation defect in the Metal compiler's inliner for large aggregate return types.
\emph{Workaround:} mandatory \texttt{\_\_attribute\_\_((noinline))} on all functions returning structs larger than $\sim$64~bytes.
\emph{Status:} not reported to Apple (no public Metal compiler bug tracker); confirmed present in Xcode~15.x.

\paragraph{Removing identity checks in bucket summation.}
Hypothesis: eliminating the identity-point check in MSM bucket accumulation (``skip if bucket is empty'') removes a branch, improving GPU warp coherence.
Result: 27\% \emph{worse}.
Root cause: the majority of buckets are empty for typical scalar distributions; skipping empty adds saves more work than the branch divergence costs.
The branch is \emph{well-predicted} by the GPU because empty buckets cluster spatially after sorting.
\emph{Lesson:} on Metal GPUs, data-dependent early-exit branches can be profitable when the skip rate is high and spatially coherent, unlike the conventional CUDA wisdom of avoiding divergent branches.

\paragraph{Uninitialized output buffers via \texttt{storageModePrivate}.}
Hypothesis: allocating output buffers as GPU-private (skipping zero-initialization) saves $\sim$3\,ms.
Result: not applicable---all buffers in a unified memory system require CPU access for result readback.
Buffer caching (grow-only allocation, reuse across calls) already eliminates repeated zeroing, making the one-time initialization cost negligible.
\emph{Lesson:} unified memory fundamentally constrains the buffer storage mode design space.
Optimizations that exploit GPU-private memory on discrete architectures do not translate.

\paragraph{MSM scalar recoding variants.}
Three alternatives to the signed-digit decomposition were tested from the 23\,ms baseline:
(a)~function inlining of the recoding loop,
(b)~segment-based partial recoding,
(c)~Horner-style evaluation of partial sums.
All three regressed.
Inlining increased register pressure beyond the compiler's ability to spill efficiently; segmentation added dispatch overhead that exceeded any per-segment savings; Horner evaluation introduced serial dependencies into what was previously a fully parallel reduction.
\emph{Lesson:} at the 23\,ms level, the MSM pipeline is tightly balanced between compute, memory, and dispatch costs.
Local improvements to any one dimension tend to regress another, consistent with Amdahl's law applied to a multi-bottleneck system.

\paragraph{Summary.}
Table~\ref{tab:dead-ends} provides a condensed catalog.
The dead ends cluster into three patterns:
(1)~\emph{Metal hardware assumptions}: 64-bit limbs, SIMD shuffles for large types, and low-occupancy kernels all fail because Metal's ALU width, shuffle granularity, and scheduler behavior differ from CUDA in ways that are undocumented;
(2)~\emph{Unified memory cost model}: local binning, private buffers, and readback elimination are optimizations designed for discrete GPU memory hierarchies that do not apply when CPU and GPU share a cache;
(3)~\emph{Algorithmic sensitivity to hardware}: GLV, lazy reduction, and fused kernels are sound algorithms whose profitability depends on hardware-specific constants (bucket count pathologies, value growth rates, barrier costs) that can only be determined empirically.

\begin{table}[t]
\centering
\small
\caption{Dead-end catalog (selected; 30+ total). $\Delta$: change from baseline; negative is regression.}
\label{tab:dead-ends}
\begin{tabular}{p{2.8cm}p{1.3cm}rp{2.2cm}}
\toprule
Optimization & Primitive & $\Delta$ & Root Cause \\
\midrule
64-bit limbs & MSM & $-$67\% & No native u64 mul \\
SIMD cooperative & MSM & $-$140\% & 96B shuffle cost \\
ARM64 asm Montgomery & CPU all & $-$23\% & LLVM already optimal \\
Fused reduce+bucket & MSM & $-$21\% & Low occupancy \\
GLV on secp256k1 & MSM & $-$12$\times$ & Scheduling pathology \\
GLV on BLS12-377 & MSM & $-$27\% & 12-limb point cost \\
Local-binning sort & Sort & $\pm$0\% & L2 cache reduces amplification \\
Lazy P2 reduction & Hash & infeasible & Value overflow at 2 rounds \\
Fused Blake3 Merkle & Hash & $-$100\% & Barrier $>$ compute \\
Transpose NTT $2^{22}$ & NTT & $-$37\% & 4 blit cost $>$ coalesce benefit \\
traceLDE no readback & STARK & $-$38$\times$ & CPU FRI needs data \\
Inline miscompilation & MSM & incorrect & Metal compiler bug \\
Remove identity check & MSM & $-$27\% & Empty-bucket skip is profitable \\
\bottomrule
\end{tabular}
\end{table}


% =============================================================================

\section{Conclusion}
\label{sec:conclusion}

We have presented zkMetal, the first comprehensive zero-knowledge cryptography library for Apple Silicon GPUs, spanning 60+~primitives across 18~fields and 10~elliptic curves.
Our contributions are threefold.

First, we identify and characterize a previously undocumented M-series GPU scheduling pathology where specific threadgroup counts cause 10--30$\times$ performance cliffs (\S\ref{sec:pathology}).
This finding has immediate practical implications for any GPU compute workload on Apple Silicon and motivates the need for better scheduling documentation from Apple.

Second, we develop a systematic methodology for field-hardware co-design on 32-bit ALU architectures (\S\ref{sec:arch}--\S\ref{sec:opt}).
The key insight is that Metal's 32-bit limitation becomes an \emph{advantage} for small fields (BabyBear NTT at 8.5\,B~elem/s, within 1.2$\times$ of the hardware floor) while requiring careful CIOS Montgomery arithmetic with explicit carry management for 256-bit fields.
The C/NEON acceleration layer achieves 10--134$\times$ over baseline Swift across all CPU primitives, demonstrating that the language-level abstraction penalty is the dominant bottleneck in CPU-side ZK computation.

Third, we provide a rigorous performance characterization against hardware-theoretical floors for 33~primitives (\S\ref{sec:floors}) and a systematic catalog of 30+~failed optimizations (\S\ref{sec:dead-ends}).
The floor analysis reveals that remaining headroom is dominated by memory access pattern irregularity (MSM scatter, NTT strided access), not compute inefficiency---a finding that redirects optimization effort toward data layout rather than arithmetic improvements.
The negative result catalog prevents duplicated effort and reveals that many CUDA-derived optimization strategies do not transfer to unified memory architectures.

Against existing implementations, zkMetal achieves 30--500$\times$ over ICICLE-Metal on NTT and 6--33$\times$ on MSM, while remaining within 3$\times$ of ICICLE-CUDA on an RTX~3090~Ti for large-field MSM---a gap explained entirely by the 6.7$\times$ memory bandwidth differential and the lack of native 64-bit integer multiply.

\paragraph{Open problems.}
The CUDA performance gap for 256-bit fields remains open: closing it likely requires either hardware changes (native 64-bit multiply in future Apple GPUs) or algorithmic innovations that restructure MSM to avoid random-access scatter entirely.
Apple's Metal compiler opacity---no PTX-equivalent IR, no public occupancy calculator, no scheduling documentation---forces empirical optimization where analytical approaches should suffice.
The inline miscompilation bug (\S\ref{sec:dead-compiler}) and scheduling pathology (\S\ref{sec:pathology}) would benefit from vendor engagement.

\paragraph{Future work.}
Three directions are most promising:
(1)~\emph{Rust FFI for ecosystem integration}: exposing zkMetal's GPU kernels to SP1, Halo2, and Plonky3 via a C-compatible FFI layer would provide immediate speedups to the most widely deployed ZK frameworks without requiring those projects to adopt Metal directly.
(2)~\emph{WebGPU port for browser proving}: Metal is the WebGPU backend on Apple devices; translating Metal shaders to WGSL would enable client-side ZK proving in web applications, addressing the growing demand for in-browser privacy-preserving computation.
(3)~\emph{Constraint IR compilation}: the runtime constraint compilation pipeline (248\,M~constraints/s, \S\ref{sec:proof-systems}) suggests that a general-purpose STARK prover accepting arbitrary AIR constraints as input and compiling them to Metal compute kernels is feasible, eliminating the need for per-circuit shader development.
